\section{Trust-Game Generative Process}
\label{app:generative_process}

This appendix specifies the environment-side process used in the repeated trust-game simulations. The generative process consists of scripted partners with hidden type and stance variables, investment-dependent stance transitions, cooperate-or-defect responses, and deterministic payoff delivery. Appendix~\ref{app:pomdp} then specifies the focal agent's POMDP model of this process.

\subsection{Episode structure}

Each episode consists of repeated interactions between one focal active-inference agent and four environment-side partners. In each round, the interaction partner is either assigned by the configured protocol or selected by the focal agent in partner-choice conditions. The focal agent then chooses a graded investment level \(a \in \mathcal{A} = \{0,1,2,3,4,5\}\), where \(a=0\) denotes no investment and \(a=5\) denotes maximum investment. The selected partner emits a binary response \(o^{\mathrm{act}}_k \in \{\mathrm{cooperate},\mathrm{defect}\}\), sampled from the partner's current type and stance. The environment then delivers payoff according to the deterministic payoff function in Appendix~\ref{app:process_payoff}, and updates the partner's stance according to the investment-dependent transition model in Appendix~\ref{app:process_stance}.

\subsection{Partner state variables}

Each environment-side partner \(k\) has two process variables hidden from the focal agent: behavioral type \(s^{\mathrm{type}}_k \in \{\mathrm{cooperator},\mathrm{reciprocator},\mathrm{exploiter},\mathrm{random}\}\) and stance toward the focal agent \(s^{\mathrm{stance}}_k \in \{\mathrm{trusting},\mathrm{neutral},\mathrm{hostile}\}\). These govern how the environment samples partner responses and updates stance over the episode; type sets the partner's broad response tendency, and stance modulates how cooperative those responses are toward the focal agent at the current interaction.

\subsection{Partner-response probabilities}
\label{app:process_response}

The partner cooperation likelihood table defines the probability the partner cooperates given type and stance:

\Needspace{14\baselineskip}
\begin{table}[H]
\centering
\caption{Cooperation likelihood $P(\text{cooperate} \mid s^{\text{type}}, s^{\text{stance}})$.}
\label{tab:coop_likelihood}
\begin{tabular}{lccc}
\toprule
Partner type & Trusting & Neutral & Hostile \\
\midrule
Cooperator & 0.95 & 0.80 & 0.55 \\
Reciprocator & 0.90 & 0.70 & 0.30 \\
Exploiter & 0.70 & 0.35 & 0.10 \\
Random & 0.60 & 0.50 & 0.35 \\
\bottomrule
\end{tabular}
\end{table}

The table defines the environment-side cooperation probability \(P(o^{\mathrm{act}}_k=\mathrm{cooperate}\mid s^{\mathrm{type}}_k,s^{\mathrm{stance}}_k)\), with \(P(o^{\mathrm{act}}_k=\mathrm{defect}\mid \cdot)=1-P(o^{\mathrm{act}}_k=\mathrm{cooperate}\mid \cdot)\).
The same response probabilities are used in the focal agent's partner-response likelihood \(A^{\mathrm{act}}_k\) (Appendix~\ref{app:pomdp_likelihoods}), allowing the focal agent to infer type and stance from observed cooperate-or-defect responses.

\subsection{Payoff function}
\label{app:process_payoff}

Payoff is delivered by the environment after the partner response is sampled. With endowment \(E=10\) and multiplier \(M=3\), the invested amount is subtracted from the focal agent's endowment. If the partner cooperates, the investment is multiplied by \(M\) and split evenly, so the focal agent receives \(Ma/2\). If the partner defects, the focal agent receives no return from the investment.
\begin{equation}
\pi(a,o_k^{\mathrm{act}})=
\begin{cases}
E-a+\frac{Ma}{2}, & o_k^{\mathrm{act}}=\mathrm{cooperate},\\
E-a, & o_k^{\mathrm{act}}=\mathrm{defect}.
\end{cases}
\label{eq:graded_payoff_map}
\end{equation}
Thus, \(a=0\) is risk-free, while higher investment increases both possible gains and exposure to defection.

\subsection{Stance transition dynamics}
\label{app:process_stance}

Stance transitions depend on the focal agent's graded investment $a$ by interpolating between a high-investment/trust-building endpoint and a low-investment/withholding endpoint:
\begin{equation}
\rho(a)=\frac{a}{|\mathcal{A}|-1}=\frac{a}{5},
\qquad
B^{\mathrm{stance}}(a)=\rho(a)B^{\mathrm{high}}+(1-\rho(a))B^{\mathrm{low}}.
\label{eq:b_stance}
\end{equation}
High investment shifts dynamics toward trust building; low investment or withholding shifts dynamics toward trust damage; recovery from hostility is slower than collapse.

\Needspace{16\baselineskip}
\begin{table}[H]
\centering
\caption{Endpoint stance transition dynamics. The focal agent's graded investment interpolates between $B^{\mathrm{high}}$ and $B^{\mathrm{low}}$ according to Eq.~\eqref{eq:b_stance}.}
\label{tab:stance_transition}
\begin{tabular}{lccc}
\toprule
\multicolumn{4}{c}{\textbf{High investment / trust-building endpoint }$B^{\mathrm{high}}$} \\
\midrule
To $\setminus$ From & Trusting & Neutral & Hostile \\
\midrule
Trusting & 0.90 & 0.30 & 0.05 \\
Neutral & 0.10 & 0.60 & 0.35 \\
Hostile & 0.00 & 0.10 & 0.60 \\
\midrule
\multicolumn{4}{c}{\textbf{Low investment / withholding endpoint }$B^{\mathrm{low}}$} \\
\midrule
To $\setminus$ From & Trusting & Neutral & Hostile \\
\midrule
Trusting & 0.10 & 0.05 & 0.02 \\
Neutral & 0.50 & 0.35 & 0.18 \\
Hostile & 0.40 & 0.60 & 0.80 \\
\bottomrule
\end{tabular}
\end{table}

These dynamics instantiate the asymmetry of trust: sustained high investment gradually builds trust, while sustained withholding rapidly damages trust; recovery from hostility is slow. In scheduled betrayal and repair protocols, the environment can additionally intervene on a partner's stance according to the protocol schedule; these interventions are not known transition events in the focal agent's POMDP.

\subsection{Boundary of the current simulation}

All reported simulations use one focal active-inference agent interacting with four scripted partners that maintain hidden type and stance, sample responses by Appendix~\ref{app:process_response}, and update stance by Appendix~\ref{app:process_stance}; they do not infer policies or maintain affective precision.
